\section{Introduction}

We study whether irrationality survives product with one projective line.
Stable rationality allows a projective-space factor of arbitrary dimension;
here the dimension of the factor matters. Let \(X\) be a smooth complex cubic
threefold. Clemens and Griffiths proved that \(X\) is irrational \cite{CG}.
We prove that its first stabilization is irrational as well.

\begin{theorem}[One-step irrationality]
\label{thm:every-cubic}
\coverage{conditional_deduction}%
\lean{rankTwoResidueMarker_eq_of_birational,%
  cubicThreefold_oneProjectiveLine_conclusion_of_residueMarker}%
For every smooth complex cubic threefold \(X\), the fourfold
\(X\times\PP^1\) is irrational.
\end{theorem}

The proof constructs an integer that is two on $X\times\PP^1$ and zero
on $\PP^4$. It counts selected two-dimensional generalized eigenspaces
of quantum multiplication. The essential step is to show that every point,
curve and surface contributes zero to this count in the quantum blowup
formula. Weak factorization then makes the count a birational invariant
of fourfolds, proving the theorem.

Surface centers explain the difficulty. In the classical blowup formula,
a surface $S$ contributes $H^3(S)(-1)$ in degree five, where
$H^3(X)(-1)$ also occurs in $H^5(X\times\PP^1)$. The classical
intermediate-Jacobian argument does not directly give the required
obstruction. The quantum selection separates the cubic
contribution from every possible center contribution.

The cubic calculation is the model case of a construction covering the
entire classification of smooth Fano threefolds of Picard rank one.
Here ``one stabilization'' always means product with $\PP^1$.

\Needspace{10\baselineskip}
\begin{theorem}[Rationality after one stabilization]
\label{thm:fano-one-step-classification}\coverage{absent}%
For every smooth complex Picard-rank-one Fano threefold $X$,
\[
 X\times\PP^1\text{ is rational}\quad\Longleftrightarrow\quad
 X\text{ is rational}.
\]
The rational families are $\PP^3$, the smooth quadric, index-two degrees
$4,5$, and index-one genera $7,9,10,12$. The other nine families have
irrational first stabilization.
\end{theorem}

Beyond detecting irrationality, we ask what the first stabilization
remembers about the original threefold. A second result recovers its
rational third cohomology for the nine families detected by the numerical
argument, denoted by $\mathcal F$.

\begin{theorem}[Rational Hodge conservation]
\label{thm:hodge-conservation}\coverage{absent}%
Let $X,Y$ be smooth complex Picard-rank-one Fano threefolds, each of index
two and degree $1,2,3$, or of index one and genus $2,3,4,5,6,8$.
If $X\times\PP^1$ and $Y\times\PP^1$ are birational, then
\[
 H^3(X,\Q)\cong H^3(Y,\Q)
\]
as rational Hodge structures.
\end{theorem}

Both theorems apply to every smooth member of the stated families.
The retained Hodge structure can distinguish birational classes among
irrational products. Combined with a separate rational Torelli theorem,
it recovers a very general cubic or quartic from its first stabilization
among smooth hypersurfaces of the same degree (Section~\ref{sec:torelli}). The Hodge isomorphism carries no assertion about
integral lattices or polarizations. The numerical theorem is independent
of this refinement.

The stable-rationality picture varies across the moduli space.  Hassett and
Tschinkel identified varieties birational to cubic threefolds as the
exceptional Fano-threefold case not covered by their stable-irrationality
theorem \cite{HassettTschinkel}.  More recently, Engel, de~Gaay Fortman, and
Schreieder proved that a very general complex cubic threefold is not stably
rational \cite[Corollary~1.4]{EGFS}, whereas Tschinkel and Zhang constructed
explicit stably rational smooth cubic threefolds over \(\mathbf Q\)
\cite[Theorem~1.3]{TschinkelZhangTorsors}. Theorem~\ref{thm:every-cubic}
applies uniformly in both situations: the first stabilization of those
stably rational examples is irrational. The companion paper
\emph{Sharpness of irrationality after one stabilization for cubic threefolds}
\cite{RuddSharpness} proves that this bound is sharp by giving two explicit
examples for which \(X\times\PP^2\) is rational.

The proof uses the quantum \(D\)-module (QDM): the cohomology module over the
ring of Novikov and bulk parameters, equipped with its flat quantum connection
\cite{IritaniQDM}.
The Novikov parameters record curve classes, while the bulk parameters are
coordinates on cohomology.  By the even QDM we mean the restriction to even
cohomology with the odd bulk variables set to zero.  We then
replace the ring of numerical Novikov and even bulk
parameters by its fraction field; this is what \emph{generic} means below.  In
particular, it refers to quantum parameters, not to a general point in the
moduli space of cubic threefolds.  In
the direction of the connection variable \(z\), the leading coefficient is
quantum multiplication by the Euler field.  Over an algebraic closure of the
generic coefficient field, its whole generalized eigenspaces, called
\emph{primary factors}, decompose the formal connection into blocks.

The count selects rank-two blocks for which the centered leading term \(N\) of
Euler multiplication is nonzero and satisfies \(N^2=0\). The lattice
modification in Section~\ref{sec:atomic-one-step} extracts two residue
eigenvalues; flatness keeps their difference fixed as the quantum
parameters vary.  We count the block when the two residue eigenvalues are
distinct modulo \(\Z\), equivalently when the two formal exponent classes are
distinct.
Beauville's quantum product formulas give exponent representatives
\(-1/6\) and \(-5/6\) for the cubic block, with difference \(2/3\).
These values already appear in Katzarkov--Lee--Svoboda--Petkov
\cite[pp.~375--376]{InterpretationsSpectra}; we give the residue calculation
in our conventions.  The small Gromov--Witten product formula and flatness in all even bulk
directions show that it doubles after product with \(\PP^1\); see
Lemma~\ref{lem:special-projective-line}. This proves the cubic theorem.

This use of Euler eigenspace blocks in birational geometry belongs to the same
circle of ideas as the Hodge-atom framework of
Katzarkov--Kontsevich--Pantev--Yu \cite{KKPYYu}.  Their framework combines the
generalized eigenspaces of Euler multiplication, Iritani's blowup theorem, and
weak factorization to produce birational invariants.  In particular, their
Example~6.21 uses the zero spectral piece, enhanced by the Serre operator and
formal monodromy, to recover the irrationality of every smooth cubic
threefold \cite[Example~6.21]{KKPYYu}.  Recent applications and refinements
include work on cubic, Gushel--Mukai, and K\"uchle fourfolds
\cite{GuereCubicFourfold,BenedettiManivelPerrin,
BenedettiFayGuereManivelPerrin}.

For the one-stabilization problem, we specialize Iritani's blowup
comparison so that the center summands remain distinct, and compute the
product with one projective line from its small quantum connection
\cite{IritaniBlowup,Behrend}.  Gyenge gives a complementary
formal tensor-product theorem for quantum \(D\)-modules of products and proves
multiplicativity for Hodge atoms \cite{GyengeProduct}.  Cai also
computes the same two exponent representatives and uses them with Iritani's
blowup theorem to prove symplectic irrationality of cubic threefolds
\cite{Cai}; that argument does not address stabilization.

Keeping distinct eigenvalues of the canonical modified residue, even when
their difference is integral, defines the stronger count $I_{\mathrm{lat}}$.
It detects index-two degree two and index-one genus six, where exponent
classes alone do not distinguish the repeated factor. In genera two through
five the repeated factor has even rank three, and its odd dimension gives
the obstruction. Cai's quartic persistence theorem is a precedent for the
cyclic-cluster argument \cite[Theorem~5.3]{CaiQuartic}.
For the Hodge refinement we use Iritani's equivariant blowup comparison
\cite{IritaniHodge}. Cavenaghi--Katzarkov--Kontsevich already use atomic
doubling under a $\PP^1$ product to obstruct equivariant linearization,
and retain curve-Jacobian isogeny classes in equivariant threefold invariants
\cite[Theorems~E, F and H]{AtomsSymbols}. Benedetti--Gu\'er\'e--Manivel--Perrin
obtain rational-Hodge restrictions on cubic and Gushel--Mukai fourfolds
birational to a Verra fourfold \cite[Theorem~1]{VerraPartners}.
Here the selected factors recover the whole rational $H^3$ for any smooth
members of the nine stated families whose first stabilizations are birational.

The proof routes are separated as follows. The numerical results use
neither Part~\ref{part:hodge} nor the optional general-bundle theorem.
\begin{center}\small
\begin{tabular}{@{}p{0.22\linewidth}p{0.70\linewidth}@{}}
\toprule
Result & Required route\\\midrule
Cubic irrationality & Numerical invariant and center vanishing
(Sections~\ref{sec:atomic-one-step}--\ref{sec:blowup-centers});
cubic calculation and $\PP^1$ product (Section~\ref{sec:cubic-proof}).\\
Fano classification & Numerical route; rank-three odd dimension and family
endpoints (Section~\ref{sec:fano-classification}); geometric matrix
identifications (Appendix~\ref{sec:fano-data}).\\
Hodge conservation & Numerical selection and endpoints; equivariant
fixed-base comparison and rational descent (Part~\ref{part:hodge}).\\
\bottomrule
\end{tabular}
\end{center}
Readers following only the cubic route can stop after
Section~\ref{sec:cubic-proof}. The arithmetic application and the
additive and general-bundle extensions are in the appendices.
The latter extensions retain Iritani--Koto's theorem.
Appendix~\ref{sec:motivic} also explains the dimensional limits and the
obstruction to extending this mechanism to a second stabilization.
